\section*{Hoofdstuk \ref{ch:g5:natural-hazards} --- Natuurgevaren}

\begin{solution}{exo:g5:natural-hazards:1}
(a) Alleen een \omterm{def:g5:natural-hazards:hazard}{natuurgevaar}: niemand wordt getroffen. (b) Een \omterm{def:g5:natural-hazards:disaster}{natuurramp}:
er worden huizen beschadigd. (c) Alleen een \omterm{def:g5:natural-hazards:hazard}{natuurgevaar}. (d) Een
\omterm{def:g5:natural-hazards:disaster}{natuurramp}.
\end{solution}

\begin{solution}{exo:g5:natural-hazards:2}
Uit de lucht: een hittegolf, een \omterm{def:g5:natural-hazards:storm}{storm}, een \omterm{def:g5:natural-hazards:drought}{droogte}. Uit het water: een
\omterm{def:g5:water-shapes-the-land:flood}{overstroming}. Uit de grond: een \omterm{def:g5:natural-hazards:landslide}{aardverschuiving}, een vulkaanuitbarsting,
een \omterm{def:g4:volcanoes-earthquakes-first-look:earthquake}{aardbeving}.
\end{solution}

\begin{solution}{exo:g5:natural-hazards:3}
Een \omterm{def:g5:natural-hazards:drought}{droogte} is een lange periode, van maanden of jaren, met veel minder
\omterm{def:g4:clouds-rain-snow:precipitation}{regen} dan normaal. Ze kan boeren zonder oogst laten, omdat de gewassen
verwelken en afsterven, en ze kan steden zonder genoeg drinkwater laten,
omdat \omterm{def:g1:landscapes-around-us:water}{rivieren} en \omterm{def:g1:landscapes-around-us:water}{meren} krimpen.
\end{solution}

\begin{solution}{exo:g5:natural-hazards:4}
Bijvoorbeeld: alles wat de wind kan meenemen naar binnen halen of
vastzetten (tuinstoelen, speelgoed, vuilnisbakken); de zaklampen
controleren en reservebatterijen en water klaarzetten; tijdens de \omterm{def:g5:natural-hazards:storm}{storm}
binnen blijven, uit de buurt van de ramen; naar de radio blijven luisteren
voor nieuws.
\end{solution}

\begin{solution}{exo:g5:natural-hazards:5}
(a) Een orkaan. (b) Een tyfoon. (c) Een tropische cycloon.
\end{solution}

\begin{solution}{exo:g5:natural-hazards:6}
Groen, geel, oranje, rood. Bij rood is er groot gevaar: kom meteen in actie
en volg de instructies van de mensen die de leiding hebben.
\end{solution}

\begin{solution}{exo:g5:natural-hazards:7}
Een deel van de heuvelhelling rechts is losgebroken en naar beneden
gegleden, en heeft een kaal litteken achtergelaten. De modder heeft zich
over de hele dalbodem verspreid en de weg bedekt. Hij heeft ook de \omterm{def:g1:landscapes-around-us:water}{rivier}
versperd als een dam: het rivierwater, dat niet meer kan wegstromen, heeft
zich erachter opgehoopt en het bos stroomopwaarts onder water gezet.
\end{solution}

\begin{solution}{exo:g5:natural-hazards:8}
Lena heeft gelijk. We kunnen de \omterm{def:g4:clouds-rain-snow:precipitation}{regen} of de \omterm{def:g1:landscapes-around-us:water}{rivier} niet tegenhouden, maar
we kunnen mensen buiten gevaar houden. Bijvoorbeeld: huizen op hoger
gelegen grond bouwen in plaats van op de \omterm{def:g5:water-shapes-the-land:flood}{overstromingsvlakte}; mensen op
tijd waarschuwen voor een \omterm{def:g5:water-shapes-the-land:flood}{overstroming}, zodat ze naar hoger gelegen grond
gaan; iedereen leren nooit door overstromingswater te lopen of te rijden.
\end{solution}

\begin{solution}{exo:g5:natural-hazards:9}
(a) $4 \times 60 = 240$ millimeter. (b) Een kwart van 60 is 15; in vier
maanden $4 \times 15 = 60$ millimeter. (c) Er ontbrak $240 - 60 = 180$
millimeter: drie kwart van de gewone \omterm{def:g4:clouds-rain-snow:precipitation}{regen}.
\end{solution}

\begin{solution}{exo:g5:natural-hazards:10}
(a) $1\,200 \div 4 = 300$ woningen. (b) Een derde van 300 is 100; dus staan
er nog $300 - 100 = 200$ woningen op de \omterm{def:g5:water-shapes-the-land:flood}{overstromingsvlakte}.
\end{solution}

\begin{solution}{exo:g5:natural-hazards:11}
Mogelijke redenen: het eerste dorp was op tijd gewaarschuwd en de mensen
hadden zich voorbereid (luiken dicht, spullen naar binnen gehaald), het
tweede niet; de huizen van het eerste zijn stevig gebouwd, die van het
tweede zijn kwetsbaar; het tweede dorp ligt aan een lage \omterm{def:g1:landscapes-around-us:water}{kust} of op een
\omterm{def:g5:water-shapes-the-land:flood}{overstromingsvlakte}, waar de \omterm{def:g5:natural-hazards:storm}{storm} ook een \omterm{def:g5:water-shapes-the-land:flood}{overstroming} bracht, terwijl het
eerste op hoger gelegen grond ligt; in het tweede dorp vielen bomen vlak bij
de huizen op de huizen. Hetzelfde gevaar richt niet overal dezelfde schade
aan: het hangt ervan af waar en hoe mensen wonen, en hoe goed ze voorbereid
zijn.
\end{solution}
